\section{Preliminaries and Formalization\iflongversion: Reward Model Robustness\fi} \label{sec:background}

Given a prompt $x$ and a response $y$, a RM produces a score $\hat{s}=RM(x,y)$.
RMs can be trained on a dataset of scored responses $D=\{(x,y,s)\}$ using (for example) a regression objective, minimizing
\begin{align} \label{eq:orig-objective}
    \E_{(x,y,s)\sim D}\left[(RM(x,y)-s)^2\right].
\end{align}
This RM training process is usually initialized from an ``SFT'' model (autoregressively pretrained LM that has been subsequently finetuned on instruction data; \citealp{bai2022training,ouyang2022training}).
Alternatively, RMs may be trained using a dataset of pairwise preferences under a Bradley-Terry assumption~\citep{bradley-terry}, maximizing
\begin{align}
    \E_{(x,y_w,y_l)\sim\mathcal{D}}[\log \sigma \left(r(x,y_w) - r(x,y_l)\right)]
\end{align}
where $y_w$/$y_l$ are the winning/losing responses.


The dataset $D$ may contain spurious correlations (e.g., with longer responses more frequently preferred; \citealp{singhal2023long}) that cause the RM to overfit to such artifacts and fail to generalize to out-of-distribution samples.
This has been observed in other classification/regression tasks~\citepia{gururangan-etal-2018-annotation,poliak-etal-2018-hypothesis,mccoy-etal-2019-right},
but is especially important for RMs. 
First, their usually small training sets (due to the cost of data collection that in principle requires human judgment) are prone to be overfit (e.g. to stylistic artifacts).
Second, RMs are expected to be robust to a wide test-time distribution: when used as evaluators, they need to judge diverse LM-generated outputs; when used in alignment, any overfitting effect would be actively \emph{exploited} by policy models, leading to ineffective alignment~\citepia{gao2022scaling,coste2024reward,eisenstein2024helping}.

We
\iflongversion
operationalize
\else
measure
\fi
RM robustness by its consistency under equivalence-maintaining transformations.
For transformed inputs
\iflongversion
$\tilde{x}, \tilde{y} = \delta(x, y)$\footnote{In special cases, some transformations also need to know if $y$ is the chosen one (see \S\ref{sec:controlled-transformations}): $\tilde{x}, \tilde{y} = \delta(x, y, \mathbb{I}[y=y_w])$.}
\else
$\tilde{x}, \tilde{y} = \delta(x, y)$
\fi
with the same meaning as the originals, an ideal RM should assign similar scores: $RM(x,y)\approx RM(\tilde{x}, \tilde{y})$.
This is a standard
\iflongversion
formalization of robustness in ML~\citepia{42503,43405,8424625}.
\else
robustness formalization~\citep{42503,43405,8424625}.
\fi
It is complementary to previous studies on RM
\iflongversion
\emph{sensitivity}, examining prediction changes
\else
\emph{sensitivity}
\fi
under meaning-\emph{altering}
\iflongversion
transformations~\citep{shen2024the}, following another line of work in NLP~\citepia{Kaushik2020Learning,gardner-etal-2020-evaluating}.
\else
transformations~\citep{shen2024the}.
\fi

\paragraph{Our focus on ranking robustness.}
It is challenging for transformations to exactly maintain equivalence. For example, wrapping the response with quotation marks maintains semantic equivalence but can be considered having worse style which would justify a lowered score.
Thus, we mainly consider the ranking that an RM assigns to a response pair, $y_w$ and $y_l$, expecting $\mathbb{I}[RM(x,y_w)>RM(x,y_l)]=\mathbb{I}[RM(\tilde{x},\tilde{y}_w)>RM(\tilde{x},\tilde{y}_l)]$ with transformed $\tilde{x}$, $\tilde{y}_w$, and $\tilde{y}_l$ where $\mathbb{I}[\cdot]$ is the indicator function.
E.g., when quotation marks are applied to both $y_w$ and $y_l$, stylistic changes equally affect both, and \emph{the RM ranking should not change}.
